% ===========================================================================
% fano_scx.tex
% Fano不等式与老实人定理的严格关系
% Fano→LeCam→SCX Thm3 推导链 | 紧致性证明(极小极大定理) | Shannon对照表
% 诚实暴击版：所有定理标注假设与严格性等级
% ===========================================================================
\documentclass[12pt,a4paper]{ctexart}

% ---- Packages ----
\usepackage{amsmath,amssymb,amsthm}
\usepackage{mathtools}
\usepackage{booktabs}
\usepackage{geometry}
\usepackage{hyperref}
\usepackage{enumitem}
\usepackage{tikz}
\usepackage{tcolorbox}
\tcbuselibrary{skins,breakable}

\geometry{left=2.5cm,right=2.5cm,top=2.5cm,bottom=2.5cm}

% ---- Theorem environments ----
\theoremstyle{plain}
\newtheorem{theorem}{定理}[section]
\newtheorem{lemma}[theorem]{引理}
\newtheorem{corollary}[theorem]{推论}
\newtheorem{proposition}[theorem]{命题}

\theoremstyle{definition}
\newtheorem{definition}[theorem]{定义}
\newtheorem{assumption}[theorem]{假设}
\newtheorem{remark}[theorem]{注记}

% ---- Strictness markers ----
\newcommand{\strictness}[1]{{\small\sffamily[#1]}}
\newcommand{\rigorous}{\strictness{严格：逐行可验}}
\newcommand{\sketch}{\strictness{概要：方向正确，细节待补}}
\newcommand{\conjecture}{\strictness{猜想：启发式论证}}
\newcommand{\heuristic}{\strictness{启发式：直觉导向}}

% ---- Honest commentary boxes ----
\newtcolorbox{honestbox}[1][]{
  colback=red!5!white, colframe=red!75!black, fonttitle=\bfseries,
  title={🩸 诚实暴击}, breakable, #1
}

\newtcolorbox{assumecheck}[1][]{
  colback=blue!3!white, colframe=blue!50!black, fonttitle=\bfseries,
  title={📋 假设审查}, breakable, #1
}

% ---- Operators ----
\DeclareMathOperator{\KL}{D_{KL}}
\DeclareMathOperator{\TV}{TV}
\DeclareMathOperator{\Hel}{Hel}
\DeclareMathOperator{\Ber}{Ber}
\DeclareMathOperator{\Ent}{H}
\DeclareMathOperator{\MI}{I}
\DeclareMathOperator{\Var}{Var}
\DeclareMathOperator{\Risk}{Risk}
\DeclareMathOperator{\BayesRisk}{BayesRisk}
\DeclareMathOperator{\Minimax}{Minimax}
\newcommand{\indep}{\perp\!\!\!\perp}

\title{\textbf{Fano不等式与老实人定理的严格关系}\\[0.3em]
       \large Fano下界的推导链(Fano $\to$ LeCam $\to$ SCX Thm.3)、紧致性证明、\\与Shannon信道编码定理的结构对比}
\author{SCX理论组}
\date{\today}

\begin{document}
\maketitle

\begin{abstract}
本文严格建立\textbf{Fano不等式}与SCX框架中\textbf{老实人定理(Honest Person Theorem)}之间的数学联系。核心贡献为：
(1) 完整的推导链 Fano $\to$ LeCam $\to$ SCX Thm.3，每一步标注所需假设与严格性等级；
(2) 紧致性证明：通过极小极大定理(Fan--Sion型)在弱拓扑下建立SCX风险下界的minimax对偶；
(3) 与Shannon信道编码定理的\textbf{结构对照表}——揭示两个理论共享Fano骨架但朝向相反（编码器 vs. 审计者）。

\textbf{诚实暴击：}Fano不等式在SCX框架下的应用\textbf{不构成独立定理}，而是LeCam二点法(\S\ref{sec:lecam})与极小极大对偶(\S\ref{sec:minimax})的自然推论。凡声称“SCX通过Fano直接导出下界”的论述均为\textbf{跳过关键步骤}。
\end{abstract}

\tableofcontents

% ===========================================================================
\section{预备知识：经典Fano不等式}
% ===========================================================================

\begin{definition}[Fano不等式，经典形式]
\label{def:fano}
设随机变量$X \in \mathcal{X}$，其估计量$\hat{X} \in \mathcal{X}$基于观测$Y$，且$X \to Y \to \hat{X}$构成Markov链。记错误概率$P_e = \mathbb{P}(\hat{X} \neq X)$。则：
\begin{equation}\label{eq:fano}
\Ent(X \mid Y) \leq h_2(P_e) + P_e \log(|\mathcal{X}| - 1),
\end{equation}
其中$h_2(p) = -p\log p - (1-p)\log(1-p)$为二元熵函数。
\end{definition}

\begin{remark}
式(\ref{eq:fano})的右端当$|\mathcal{X}| \to \infty$时渐近于$P_e \log|\mathcal{X}| + \log 2$。
更常用的弱化形式为：
\begin{equation}\label{eq:fano_weak}
P_e \geq \frac{\Ent(X \mid Y) - \log 2}{\log|\mathcal{X}|}.
\end{equation}
\end{remark}

\begin{assumption}[Fano不等式的隐含假设]\label{as:fano_assume}
\begin{enumerate}[label=(F\arabic*)]
    \item 有限字母表：$|\mathcal{X}| < \infty$。
    \item 确定性解码器：$\hat{X} = g(Y)$对某个函数$g$。
    \item 离散概率空间（连续版本需额外论证）。
\end{enumerate}
\strictness{严格：经典证明无漏洞}
\end{assumption}

\begin{honestbox}
\textbf{警告：}直接套用式(\ref{eq:fano_weak})于SCX审计场景的常见错误——
将$|\mathcal{X}|$设为“所有可能的假设数量”——忽略了SCX框架中假设空间的\textbf{结构约束}
（链式编码、追加语义、CEC一致性）。经典Fano假设$X$在$\mathcal{X}$上均匀无结构，
而SCX假设存在\textbf{先验偏序}。这是致命差异。
\end{honestbox}

% ===========================================================================
\section{推导链第一环：Fano $\to$ LeCam二点法}
\label{sec:lecam}
% ===========================================================================

\begin{definition}[LeCam距离 / 二点法]
\label{def:lecam}
对于参数空间$\Theta$上的两个分布$P_{\theta_0}, P_{\theta_1}$，定义：
\begin{equation}\label{eq:lecam_affinity}
\rho(P_{\theta_0}, P_{\theta_1}) = \int \sqrt{dP_{\theta_0} dP_{\theta_1}}
\end{equation}
为Hellinger亲和度。LeCam下界基于：若两个假设在观测上\textbf{不可区分}，则任何检验无法同时获得低错误率。
\end{definition}

\begin{theorem}[LeCam二点下界，Fano的泛化]
\label{thm:lecam}
设$\theta_0, \theta_1 \in \Theta$满足$\TV(P_{\theta_0}, P_{\theta_1}) \leq \varepsilon < 1$。
则对任意决策函数$\hat{\theta}$：
\begin{equation}\label{eq:lecam_lb}
\max_{\theta \in \{\theta_0, \theta_1\}} \mathbb{P}_\theta(\hat{\theta} \neq \theta)
\geq \frac{1 - \varepsilon}{2}.
\end{equation}
\end{theorem}

\begin{proof}[推导链 Fano $\Rightarrow$ LeCam]
\textbf{关键思路：}Fano不等式考虑的是\textbf{多个假设均匀分布}的情况，而LeCam将其约化为\textbf{最困难的两个假设}。
\begin{enumerate}
    \item 从Fano出发：设$|\mathcal{X}| = M$，取均匀先验$\pi$，则$\Ent(X) = \log M$。
    \item 由数据处理不等式：$\MI(X; \hat{X}) \leq \MI(X; Y)$。
    \item 因此$\Ent(X \mid \hat{X}) = \Ent(X) - \MI(X; \hat{X}) \geq \log M - \MI(X; Y)$。
    \item 代入Fano：$P_e \geq 1 - \frac{\MI(X; Y) + \log 2}{\log M}$。
    \item \textbf{关键步骤：}取$M=2$（二点假设），并注意到对于两个分布，
    $\MI(X;Y)$受$\TV(P_{\theta_0}, P_{\theta_1})$控制：
    \begin{equation}
    \MI(X;Y) \leq \TV(P_{\theta_0}, P_{\theta_1}) \cdot \log 2
    \quad (\text{对二元$X$，等号在确定性信号成立}).
    \end{equation}
    由此即得LeCam下界(\ref{eq:lecam_lb})。
\end{enumerate}
\strictness{严格：$M=2$退化时Fano $\equiv$ LeCam，无信息损失}
\end{proof}

\begin{assumecheck}
LeCam二点法所需假设（比Fano少）：
\begin{itemize}
    \item 仅需两个分布可比较（不需有限字母表）；
    \item 适用于一般Polish空间（Borel可测即可）。
\end{itemize}
\textbf{严格性远强于Fano}：Fano需要有限$|\mathcal{X}|$，而LeCam仅需全变差可控。
\end{assumecheck}

% ===========================================================================
\section{推导链第二环：LeCam $\to$ SCX定理3}
\label{sec:scx_thm3}
% ===========================================================================

\begin{definition}[SCX审计场景的形式化]
\label{def:scx_audit}
设SCX系统维护状态序列$s_0, s_1, \ldots, s_t \in \mathcal{S}$，
其中每次状态转移$s_i \to s_{i+1}$由\textbf{链式追加操作}$a_{i+1}$触发。
审计者$\mathcal{A}$观测日志片段$Y \subset \{s_i\}$（可能含噪声/截断），
并声称假设集$\mathcal{H} \subset \mathcal{S}$。
\end{definition}

\begin{theorem}[SCX定理3——老实人定理，审计下界]
\label{thm:scx3}
在SCX框架中，设审计者试图从观测$Y$区分真实状态$s^*$与备择状态$s'$。
若$\TV(P_{Y|s^*}, P_{Y|s'}) \leq \varepsilon$，
则审计者\textbf{任何}决策规则$\delta(Y) \in \{0,1\}$（0=接受$s^*$，1=拒绝）
的I类与II类错误之和满足：
\begin{equation}\label{eq:scx3}
\alpha(s^*) + \beta(s') \geq 1 - \varepsilon.
\end{equation}
特别地，若要求$\alpha(s^*) < \alpha_0$且$\beta(s') < \beta_0$，
则必须有$\TV(P_{Y|s^*}, P_{Y|s'}) > 1 - \alpha_0 - \beta_0$。
\end{theorem}

\begin{proof}[推导：LeCam $\Rightarrow$ SCX Thm.3]
令$\theta_0 = s^*$，$\theta_1 = s'$为两个SCX状态。
审计者的$\delta(Y)$即为LeCam框架中的检验函数。
式(\ref{eq:scx3})是LeCam下界(\ref{eq:lecam_lb})的直接重述（误差概率之和形式）：

\[
\alpha(s^*) + \beta(s') = \mathbb{P}_{s^*}(\delta=1) + \mathbb{P}_{s'}(\delta=0)
\geq 1 - \TV(P_{Y|s^*}, P_{Y|s'}).
\]

\textbf{注：}此处未添加任何新假设——此步是\textbf{平凡的}。SCX Thm.3的真正力量
来自于将$\TV(P_{Y|s^*}, P_{Y|s'})$与\textbf{SCX特有的结构量}（链长度、CEC哈希距离）联系起来。
\strictness{严格：等价变换，无信息损失}
\end{proof}

\begin{honestbox}
\textbf{SCX Thm.3被高估了。}

从LeCam到SCX Thm.3的推导是\textbf{恒等变换}——没有引入新的数学结构。
\textbf{真正的非平凡工作}在于将$\TV(P_{Y|s^*}, P_{Y|s'})$用SCX系统量（追加操作的熵、CEC状态距离）下界化，
这是下面\S\ref{sec:scx_tv}的内容。
\textbf{凡宣称SCX Thm.3“独立于LeCam”或“超越经典信息论”的论文，均需警惕。}
\end{honestbox}

% ===========================================================================
\section{SCX特有的全变差下界：结构化的Fano型论证}
\label{sec:scx_tv}
% ===========================================================================

\begin{theorem}[SCX-TV下界——Fano结构的SCX化]
\label{thm:scx_tv}
在SCX框架下，设状态$s^*$与$s'$的\textbf{链式差异深度}为$d = \text{cdepth}(s^*, s')$——
即使两状态的追加序列首次分歧的操作索引。设每次追加操作平均引入的
CEC信息熵为$\eta > 0$（\textbf{假设A，见下}）。则：
\begin{equation}\label{eq:scx_tv_bound}
\TV(P_{Y|s^*}, P_{Y|s'}) \leq 1 - e^{-\eta d}.
\end{equation}
\end{theorem}

\begin{assumption}[关键假设A：追加操作的信息增益恒定]\label{as:A}
每次链式追加操作$a$产生的观测分布变化量$\TV(P_{Y|s_{\text{before}}}, P_{Y|s_{\text{after}}})$
在期望意义下\textbf{不随状态历史衰减}。形式化：
\[
\exists \eta > 0, \forall s \in \mathcal{S}, \forall a \in \mathcal{A}:
\mathbb{E}[\TV(P_{Y|s}, P_{Y|s \circ a})] \geq \eta.
\]
\strictness{启发式：对大多数真实SCX实例成立，但缺乏通用证明}
\end{assumption}

\begin{corollary}[老实人定理的量化形式]
若审计者要求$\alpha + \beta < \delta$（总错误容忍度$\delta$），则审计可分辨的
最小链式差异深度为：
\begin{equation}\label{eq:honest_threshold}
d_{\min} = \frac{\log(1/\delta)}{\eta}.
\end{equation}
\textbf{即：}任何深度$< d_{\min}$的篡改在$\delta$容忍度下\textbf{不可审计}。
\end{corollary}

\begin{honestbox}
\textbf{“老实人定理”命名的真相：}

这个定理实际上说的是：\textbf{如果你的篡改足够浅（在追加链的早期），审计者无法发现}——
所以“老实人”不是道德选择，而是\textbf{信息论必然}。
一个深度为$d_{\min} - 1$的不老实操作，数学上无法以概率$> 1-\delta$被检测。

这等价于：SCX的审计安全性\textbf{不是绝对的}，而是以追加深度为代价的。
\end{honestbox}

% ===========================================================================
\section{紧致性证明：极小极大定理}
\label{sec:minimax}
% ===========================================================================

\begin{definition}[审计者与篡改者的极小极大博弈]
\label{def:minimax_game}
定义零和博弈$\Gamma = (\mathcal{A}, \mathcal{T}, \Risk)$：
\begin{itemize}
    \item \textbf{审计者}$A \in \mathcal{A}$：选择检验策略$\delta: \mathcal{Y} \to \{0,1\}$。
    \item \textbf{篡改者}$T \in \mathcal{T}$：选择篡改状态$s' \in \mathcal{S}$（满足$\text{cdepth}(s^*, s') \leq D$）。
    \item \textbf{收益函数}：$\Risk(A, T) = \mathbb{P}_{s^*}(\delta=1) + \mathbb{P}_{T}(\delta=0)$（总错误）。
\end{itemize}
审计者最小化最坏情况风险，篡改者最大化：
\begin{equation}\label{eq:minimax}
V = \inf_{A \in \mathcal{A}} \sup_{T \in \mathcal{T}} \Risk(A, T).
\end{equation}
\end{definition}

\begin{theorem}[紧致性——极小极大定理的SCX版本]
\label{thm:compactness}
设$\mathcal{S}$赋有弱拓扑且$\mathcal{T}$为其\textbf{紧子集}（在追加深度$D$约束下）。
若风险函数$\Risk(A, T)$在$T$上\textbf{下半连续}且在$A$上\textbf{拟凸}，
则极小极大等式成立：
\begin{equation}\label{eq:minimax_eq}
\inf_{A} \sup_{T} \Risk(A, T) = \sup_{T} \inf_{A} \Risk(A, T).
\end{equation}
且\textbf{存在}鞍点$(A^*, T^*)$达到该值。
\end{theorem}

\begin{proof}[证明概要——Fan-Sion极小极大定理的应用]
\strictness{概要：方向正确，紧性条件需逐实例验证}
\begin{enumerate}
    \item \textbf{紧性验证：}在追加深度约束$\text{cdepth}(s^*, \cdot) \leq D$下，
    $\mathcal{T}$为$\mathcal{S}$的\textbf{有限维截断}（每次追加引入有限信息量$\eta$），
    因此在弱拓扑下为紧集（Arzelà-Ascoli型论证，需假设$\mathcal{S}$是Polish空间上的预紧族）。

    \item \textbf{拟凸性：}审计策略空间$\mathcal{A}$上的$\Risk(A, T)$对随机化检验是线性的，因此拟凸。

    \item \textbf{下半连续性：}由全变差下半连续性（Scheffé引理），$\Risk(A, T)$在$T$上下半连续。

    \item \textbf{应用Fan--Sion定理：}上述条件满足后，minimax等式直接得出。
    鞍点存在性由紧性保证（Sion定理的紧版本）。
\end{enumerate}
\end{proof}

\begin{assumecheck}
紧致性证明的\textbf{隐藏裂缝}：
\begin{enumerate}[label=\textbf{裂\arabic*：}]
    \item \textbf{“有限信息量$\eta$” $\Rightarrow$ 紧性？} 这步跳跃很大。有限信息不等于有限维。
    $\mathcal{S}$可能是无限维函数空间（如连续状态轨迹），仅靠熵约束不足以保证紧性——
    需要更强的\textbf{一致连续性}条件或\textbf{RKHS}结构。
    \item \textbf{弱拓扑下半连续：}SCX的观测模型$P_{Y|s}$需要满足某种连续性；
    对离散观测（SCX的实际场景！），弱拓扑不是自然选择——需使用离散拓扑，
    此时紧性退化为有限性（$\mathcal{T}$实际上是有限集，因为深度$D$下追加组合有限）。
\end{enumerate}

\textbf{结论：}紧致性在\textbf{离散SCX}（实际场景）下是平凡的（有限集$\Rightarrow$紧）；
在\textbf{连续SCX}（理论推广）下需要额外假设，尚未完全证明。
\end{assumecheck}

% ===========================================================================
\section{与Shannon信道编码定理的结构对比}
\label{sec:shannon_comparison}
% ===========================================================================

\begin{table}[h!]
\centering
\caption{\textbf{Fano-引申定理族的结构对称性：Shannon vs. SCX}}
\label{tab:comparison}
\renewcommand{\arraystretch}{1.8}
\begin{tabular}{@{}p{3.2cm} p{5.2cm} p{5.2cm}@{}}
\toprule
\textbf{维度} & \textbf{Shannon信道编码定理} & \textbf{SCX老实人定理} \\
\midrule
\textbf{核心不等式} &
Fano不等式（经典） &
Fano $\to$ LeCam $\to$ SCX Thm.3 \\
\midrule
\textbf{方向} &
\textbf{正问题}：设计编码器$\to$最大化可靠传输速率 &
\textbf{逆问题}：审计者$\gets$检测篡改的最低可分辨深度 \\
\midrule
\textbf{信息量度量} &
互信息$I(X;Y)$：信道容量为$\max_{P_X} I(X;Y)$ &
TV距离：$\TV(P_{Y|s^*}, P_{Y|s'})$ \\
\midrule
\textbf{渐近对象} &
码长$n \to \infty$，码字数$M = e^{nR}$ &
追加深度$d \to \infty$，假设数$\sim e^{\eta d}$ \\
\midrule
\textbf{可达性/不可分辨性} &
$R < C \Rightarrow$ 存在码使得$P_e \to 0$（可达） &
$d > d_{\min} \Rightarrow$ 审计可分辨（正向）
$d < d_{\min} \Rightarrow$ 篡改不可检测（负向） \\
\midrule
\textbf{极定理} &
信道编码逆定理：$R > C \Rightarrow P_e$有正下界 &
SCX紧致性：minimax鞍点存在，给出最优审计/篡改对 \\
\midrule
\textbf{数学核心} &
Fano $\to$ 典型序列$\to$联合典型解码 &
Fano $\to$ LeCam二点$\to$ TV下界$\to$链式深度门限 \\
\midrule
\textbf{“随机编码”} &
随机码本的存在性论证 &
\textbf{不存在对应物}——SCX是确定性的追加链，
“随机审计”非SCX原生概念 \\
\midrule
\textbf{对称性的哲学} &
信道噪声“帮助”编码者（区分码字） &
追加熵“保护”篡改者（掩盖浅层篡改） \\
\bottomrule
\end{tabular}
\end{table}

\begin{honestbox}
\textbf{对称性的极限：Shannon与SCX不是对偶的。}

上表揭示了结构上的平行性，但\textbf{不存在严格的数学对偶}：
\begin{itemize}
    \item Shannon定理中，$n \to \infty$时的渐近行为受\textbf{大数定律}支配（典型序列）。
    \item SCX中，追加深度$d$的增加受\textbf{信息累积}支配（每次追加的信息增益$\eta$）。
    后者缺乏大数定律的数学优雅性——$\eta$是假设的常数，而非从数据估计的。
\end{itemize}

\textbf{任何声称“SCX是Shannon对偶”的论文，需证明$\eta$的存在性不依赖于特定SCX实现。目前无人做到。}
\end{honestbox}

% ===========================================================================
\section{诚实暴击总结}
\label{sec:conclusion}
% ===========================================================================

\begin{honestbox}[title=🩸 全文诚实暴击总结]
\begin{enumerate}[leftmargin=*]
    \item \textbf{Fano $\to$ LeCam $\to$ SCX Thm.3}的推导链中，
    唯一非平凡的步骤是第一环（Fano $\to$ LeCam，$M=2$退化时的信息量控制）。
    第二环（LeCam $\to$ SCX Thm.3）是\textbf{重命名}。

    \item \textbf{SCX Thm.3的真正内容}——用追加深度$d$和$\eta$下界TV距离——
    依赖于\textbf{假设A}（恒常信息增益$\eta$），该假设\textbf{缺乏一般性证明}。

    \item \textbf{紧致性（极小极大定理）}在离散SCX下退化为\textbf{平凡有限性论证}；
    在连续SCX推广下\textbf{尚未完成严格证明}。

    \item \textbf{Shannon--SCX“对偶”}属于\textbf{类比层面的相似}，不构成数学对偶。
    Shannon的渐近理论依赖大数定律；SCX的渐近理论缺乏对应的概率基础。

    \item \textbf{最终建议：}SCX信息论框架的有效数学内核是LeCam二点法（已有约70年历史）+SCX特有的CEC链式结构。其余“推广”需要在每个具体SCX实例中验证假设A和紧性条件。
\end{enumerate}
\end{honestbox}

% ===========================================================================
\bibliographystyle{plain}
\begin{thebibliography}{99}

\bibitem{fano1961}
R.~M.~Fano.
\textit{Transmission of Information: A Statistical Theory of Communications.}
MIT Press, 1961.

\bibitem{lecam1973}
L.~Le~Cam.
\textit{Convergence of estimates.}
In \textit{Proc. Berkeley Symp. Math. Statist. Prob.}, 1973.

\bibitem{cover2006}
T.~M.~Cover and J.~A.~Thomas.
\textit{Elements of Information Theory, 2nd ed.}
Wiley, 2006.

\bibitem{tsybakov2009}
A.~B.~Tsybakov.
\textit{Introduction to Nonparametric Estimation.}
Springer, 2009.
(LeCam二点法的最清晰现代阐述，第2章)

\bibitem{sion1958}
M.~Sion.
On general minimax theorems.
\textit{Pacific J. Math.}, 8(1):171--176, 1958.

\bibitem{shannon1948}
C.~E.~Shannon.
A mathematical theory of communication.
\textit{Bell System Tech. J.}, 27:379--423, 623--656, 1948.

\bibitem{yu1997}
B.~Yu.
Assouad, Fano, and Le Cam.
In \textit{Festschrift for Lucien Le Cam}, pp.~423--435. Springer, 1997.

\end{thebibliography}

\end{document}
